\chapter{Pendahuluan}
\label{ch:01-pendahuluan}

\section{Latar Belakang dan Rumusan Masalah}

Pemberian pakan pada kolam akuaponik adalah pekerjaan harian yang paling mudah
salah: waktunya bergeser mengikuti kesibukan pemilik, dan jumlahnya ditakar dengan
perkiraan. Kedua kesalahan itu berbiaya langsung pada kualitas air. Pakan yang
berlebih tidak habis dimakan ikan, mengendap, dan terurai menjadi amonia serta
nitrit yang menekan daya dukung akuaponik; pakan yang kurang menghambat
pertumbuhan ikan dan menurunkan pasokan nutrisi bagi tanaman. Persoalan yang harus
diselesaikan karena itu bukan sekadar menyalakan motor pada jam tertentu,
melainkan menjaga pemberian pakan tetap tepat waktu dan tepat jumlah, sekaligus
memberi tahu pemilik ketika persediaan pakan menipis atau habis.

Pemantauan manual menambah persoalan kedua: pemilik tidak selalu berada di lokasi,
sehingga ia tidak tahu sisa pakan tanpa melihat langsung dan tidak dapat memberi
pakan darurat saat dibutuhkan. Rancangan ini menjawab keduanya dengan satu node
perangkat dan satu platform pemantauan. Node berbasis \BoardName{} mengukur jarak
permukaan pakan di hopper dengan sensor ultrasonik, menggerakkan gate pakan dengan
motor servo, dan menjalankan jadwal harian secara mandiri. Platform berjalan
sebagai satu proses Python yang memuat broker MQTT, worker ingesti, sweeper, basis
data SQLite, dan dashboard Streamlit.

Empat pertanyaan berikut menjadi rumusan masalah rancangan ini.
\begin{itemize}[leftmargin=*,itemsep=2pt]
  \item Bagaimana menakar pakan dalam satuan porsi yang konsisten dan melaporkan
        sisa pakan hopper sebagai persentase yang dapat dibandingkan antar waktu?
  \item Bagaimana merancang pertukaran pesan sehingga setiap perintah memiliki
        hasil yang pasti: berhasil, gagal, atau tidak dijawab sama sekali?
  \item Bagaimana menyimpan telemetri, jadwal, perintah, riwayat pakan, dan
        notifikasi dalam satu basis data yang konsisten meskipun ditulis beberapa
        thread sekaligus?
  \item Bagaimana menyajikan keadaan sistem dan sarana kendalinya dalam satu
        dashboard yang tetap ringan pada perangkat kelas laptop?
\end{itemize}

\section{Tujuan}

Sasaran perancangan ini adalah menghasilkan rancangan yang terverifikasi terhadap
kode yang benar-benar berjalan:
\begin{enumerate}[leftmargin=*,itemsep=2pt]
  \item node perangkat berbasis \BoardName{} yang menakar pakan dalam satuan
        porsi, mengukur jarak permukaan pakan secara periodik, dan tetap berfungsi
        ketika sebagian periferal tidak terpasang;
  \item kontrak pesan MQTT yang eksplisit, dengan kewajiban ack untuk setiap
        perintah dan penandaan \code{timeout} setelah \CmdTimeoutS{} detik;
  \item platform terpadu dalam satu proses Python, tanpa layanan eksternal yang
        harus dipasang lebih dahulu;
  \item skema basis data relasional berisi \NumTables{} tabel dan satu view agregat;
  \item dua jalur penjadwalan yang saling melengkapi: jadwal lokal di perangkat
        yang tetap berjalan saat jaringan putus, dan jadwal di platform yang dapat
        diubah pemilik kapan saja lalu dikirim ke perangkat.
\end{enumerate}

\section{Ruang Lingkup dan Batasan}

Rancangan mencakup perangkat keras node beserta pembagian pin dan catatan dayanya,
firmware \FwVersion{} pada folder \file{firmware/esp32_feeder}, platform pada folder
\file{dashboard} (broker tertanam, worker, sweeper, lapisan basis data, dan
\NumPages{} halaman dashboard), skema basis data pada \file{database/schema.sql},
serta kontrak protokol MQTT yang mengikat keduanya.

Batasan berikut diambil secara sadar dan berlaku untuk seluruh dokumen.
\begin{itemize}[leftmargin=*,itemsep=2pt]
  \item \textbf{Satu perangkat.} Baris perangkat awal adalah \code{\SeedDevice}
        (\SeedDeviceName), sedangkan firmware menurunkan identitas dari MAC chip
        sehingga papan kedua muncul sebagai perangkat baru, misalnya
        \code{\DeviceExample}. Perilaku banyak perangkat belum diuji dan tidak
        diklaim.
  \item \textbf{Broker anonim tanpa TLS.} Broker tertanam menerima koneksi tanpa
        autentikasi dan tanpa enkripsi, sehingga rancangan ini hanya layak untuk
        satu jaringan lokal tepercaya (\chapref{ch:09-catatan-asumsi}).
  \item \textbf{SQLite tertanam.} Basis data adalah berkas tunggal \code{\DbFile}
        dalam mode WAL; tidak ada server basis data terpisah dan tidak ada
        replikasi.
  \item \textbf{Jaringan lokal tunggal.} Alamat host dapat berubah karena DHCP,
        sehingga alamat broker diisi pada berkas \file{secrets.h} yang tidak ikut
        disimpan di repositori.
  \item \textbf{Bukan penelitian biologis.} Kualitas air dan pertumbuhan ikan hanya
        dipakai sebagai alasan perancangan; keduanya tidak diukur maupun diuji.
\end{itemize}

\section{Daftar Istilah}

Istilah teknis dipertahankan dalam bahasa Inggris karena merupakan nama konsep pada
protokol dan kode. \tabref{tab:01-istilah} merangkum istilah yang dipakai berulang.

\begin{table}[H]
\centering
\caption{Istilah yang dipakai berulang dalam dokumen ini.}
\label{tab:01-istilah}
\small
\begin{tabular}{@{}p{2.7cm}p{11.9cm}@{}}
\toprule
\textbf{Istilah} & \textbf{Arti dalam sistem ini} \\
\midrule
MQTT & Protokol publish/subscribe yang menghubungkan perangkat dengan platform. \\
Broker & Server MQTT yang meneruskan pesan antar client; di sini berjalan tertanam
di dalam proses Python aplikasi. \\
Retained & Pesan terakhir pada sebuah topik yang disimpan broker dan langsung
dikirim ke subscriber baru; dipakai hanya untuk topik status. \\
LWT & \emph{Last Will and Testament}: pesan \code{offline} yang didaftarkan
perangkat saat connect dan diterbitkan broker ketika koneksi terputus tanpa
penutupan sesi yang rapi. \\
topic & Jalur penamaan pesan berbentuk \code{<prefix>/<device_id>/<kind>}, misalnya
\topic{telemetry}. \\
QoS & Tingkat jaminan pengiriman MQTT. Arah perintah memakai QoS~1; arah telemetri
memakai QoS~0. \\
ack & Balasan perangkat atas satu perintah, memuat \code{command_id},
\code{result}, dan \code{message}; wajib untuk setiap perintah. \\
\code{command_id} & Pengenal unik perintah berbentuk \code{cmd-<uuid>} yang di-echo
perangkat dan menjadi kunci korelasi baris perintah dengan baris riwayat. \\
hopper & Wadah penampung pakan pada node, diukur ketinggian permukaannya dari atas. \\
portion & Satuan porsi: satu kali gate terbuka selama \PortionOpenMs{} ms. Satu
perintah dibatasi 1 sampai 10 porsi. \\
kalibrasi & Dua jarak acuan sensor, \HopperFull{} cm saat penuh dan \HopperEmpty{}
cm saat kosong, untuk mengubah jarak menjadi persentase. \\
WAL & \emph{Write-Ahead Logging}, mode jurnal SQLite yang memungkinkan pembaca dan
penulis bekerja bersamaan pada satu berkas. \\
sweeper & Thread berkala yang menandai perintah tanpa ack sebagai \code{timeout}
dan perangkat yang berhenti mengirim data sebagai offline. \\
\bottomrule
\end{tabular}
\end{table}

\section{Ringkasan Arsitektur dan Peta Bab}

Sistem terdiri atas dua proses yang bertemu di satu broker. Proses pertama adalah
node perangkat: satu sketch Arduino berbentuk super-loop tanpa \code{delay()}, yang
membaca sensor tiap \SensorMs{} ms, menggerakkan servo sebagai state machine,
menjalankan jadwal harian dari clock lokal, dan mempublikasikan telemetri tiap
\TelemetryMs{} ms ke topik \topic{telemetry}. Proses kedua adalah platform: satu
proses Python yang menjalankan broker MQTT tertanam pada \code{\BrokerBind}, worker
paho yang berlangganan empat jenis topik, sweeper tiap \SweepS{} detik, basis data
SQLite, dan dashboard Streamlit. Perintah mengalir balik melalui topik
\topic{command} dengan QoS~1, dan setiap perintah dijawab pada topik \topic{ack}.
Broker tertanam dipilih agar platform tidak bergantung pada layanan MQTT eksternal;
bila port sudah dipakai broker lain, platform memakai broker tersebut.

\tabref{tab:01-peta-bab} memetakan isi dokumen. Bab 2 sampai Bab 5 membahas
subsistem, Bab 6 dan Bab 7 membahas dua kontrak yang mengikatnya, Bab 8 menyatukan
keduanya beserta matriks ketertelusuran, dan Bab 9 mencatat asumsi serta temuan.

\begin{table}[H]
\centering
\caption{Peta bab dan fokus pembahasannya.}
\label{tab:01-peta-bab}
\small
\begin{tabular}{@{}p{1.0cm}p{4.0cm}p{9.6cm}@{}}
\toprule
\textbf{Bab} & \textbf{Judul} & \textbf{Fokus} \\
\midrule
\chapref{ch:02-arsitektur} & Arsitektur Sistem & Pembagian proses, komponen
runtime, pembagian thread, aliran data, pemilihan mode broker. \\
\chapref{ch:03-perangkat-keras} & Perangkat Keras & Blok perangkat keras, daftar
komponen, skema rangkaian, pembagian pin, timing servo, kalibrasi hopper. \\
\chapref{ch:04-firmware} & Firmware ESP32 & \code{setup()} dan tick kooperatif,
state machine gate, NVS, penjadwalan lokal, siklus MQTT, mode degradasi. \\
\chapref{ch:05-website} & Dashboard Website & Enam halaman, cache baca versus jalur
tulis, state sesi, internasionalisasi, sweeper, backup dan restore. \\
\chapref{ch:06-protokol-mqtt} & Protokol MQTT & Tabel topik, pasangan perintah dan
ack, contoh payload, aturan ack wajib dan \code{timeout}. \\
\chapref{ch:07-basis-data} & Basis Data & Diagram relasi, kamus data, index, view
agregat, data awal, kardinalitas, kolom yang tidak dipakai. \\
\chapref{ch:08-end-to-end} & Alur End-to-End & Perjalanan satu perintah dari klik
sampai baris basis data berubah, matriks aksi, ketertelusuran. \\
\chapref{ch:09-catatan-asumsi} & Catatan Asumsi dan Temuan & Empat temuan sisa,
asumsi lingkungan, batasan keamanan. \\
\bottomrule
\end{tabular}
\end{table}
